\frametitle{The Same Four Questions, Two Pipelines}
\small

The two pipelines are not two different subjects. They are two ways of
answering the same four questions.

\vspace{0.4em}
{\centering
\renewcommand{\arraystretch}{1.35}
\begin{tabular}{@{}p{4.6cm} p{9.4cm} p{9.4cm}@{}}
\rowcolor{black!6}
\textbf{The question} & \textbf{Classical pipeline} & \textbf{Neural-operator pipeline}\\
\hline
Can I solve it?
  & Sparse linear system, solved iteratively
  & One forward pass of the trained operator $\mathcal{G}_\theta$\\
Can I measure the error?
  & \alert{A posteriori} estimator, computable cell by cell
  & Training and validation loss, measured once, offline\\
Can I speed it up?
  & \alert{Optimal complexity} by multigrid; adaptive refinement
  & Amortised --- no solve is performed at inference time\\
Have I used every effort?
  & The smoother is the \alert{least automatic} component
  & \alert{Learn the smoother}; keep the multilevel structure\\
\end{tabular}\par}

\vspace{0.4em}
\begin{block}{What the comparison is for}
The learned route is not asked to replace the classical one. It is asked to
answer \alert{one} of the four questions --- the last --- better than a fixed
smoother can, while the other three keep their classical answers.
\end{block}
